\renewcommand{\Anno}{2001}

\renewcommand{\NumeroProblema}{4}

\renewcommand{\Testo}{Un dado perfettamente equilibrato con facce
numerate da 1 a 6 viene lanciato più volte.
\begin{enumerate}
	\item Qual è la probabilità che lanciandolo 3 volte la somma dei punteggi ottenuti sia divi\-si\-bile per 6?
	
	\item Qual è la probabilità che lanciandolo 3 volte la somma dei punteggi ottenuti sia divi\-si\-bile per 7?
	
	\item Qual è la probabilità che lanciandolo $n$ volte la somma dei punteggi ottenuti sia divi\-si\-bile per 6?
	
	\item Qual è la probabilità che lanciandolo $n$ volte la somma dei punteggi ottenuti sia divi\-si\-bile per 7?
\end{enumerate}}

\renewcommand{\Risposta}{\begin{enumerate}
		\item  $1/6$
	
		\item  $5/36$
	
		\item  Per ogni valore di $n$ la probabilità è sempre $1/6$.
	
		\item  La probabilità è data, in funzione di $n$, dalla formula
		$$\frac{1}{7}+\frac{(-1)^{n}}{7\cdot 6^{n-1}}.$$
	\end{enumerate}}

\renewcommand{\Aiutino}{Nei casi in cui siamo interessati alla
divisibilità per 6 decide sempre l'ultimo tiro!

Nei casi in cui siamo interessati alla divisibilità per 7 decide sempre l'ultimo tiro, a condizione che al penultimo la somma non fosse già multipla di 7.  Detta $p_{n}$ la probabilità che la somma sia divisibile per 7 dopo $n$ lanci, si ha quindi la ricorrenza $p_{n+1}=(1-p_{n})/6$, da cui si ottiene per induzione una formula esplicita per $p_{n}$.

Il caso $n=3$ si può ovviamente anche trattare ``a mano'',
determinando quanti sono i modi di ottenere 7 o 14 lanciando 3 dadi.}

\renewcommand{\SoluzioneUfficiale}{
}

\renewcommand{\SoluzioniAlternative}{
}

\renewcommand{\Osservazioni}{
}

\renewcommand{\Commenti}{
}